\documentclass[11pt,twocolumn]{article}
\usepackage[letterpaper,margin=0.75in]{geometry}
\usepackage[utf8]{inputenc}
\usepackage{graphicx,hyperref,booktabs,amsmath,amssymb,xcolor,float,multirow,array}
\hypersetup{colorlinks=true,linkcolor=blue,citecolor=blue}

% ── Metadata ──
\title{Where Does Scoring Bias Come From?\\A Base vs Instruct Comparison of LLM-as-a-Judge}
\author{Author Name$^{1}$ \\
$^{1}$Department, School \\
\texttt{email@school.edu}}
\date{\today}

\begin{document}
\maketitle

% ── ABSTRACT ──
\begin{abstract}
LLMs are increasingly deployed as automated judges in the LLM-as-a-Judge paradigm~\cite{zheng2023judging}, yet they exhibit systematic scoring biases that compromise evaluation reliability. Li~et~al.~\cite{li2025scoring} documented three scoring bias types across five frontier models, but explicitly noted that ``the underlying causes of these scoring biases remain to be validated.'' We present the first systematic investigation of whether scoring bias originates from pre-training or emerges during instruction tuning. Through controlled experiments spanning 30 model variants from 15 families across 3 scoring bias probes and 3 variants per probe totaling over 40,000 judgments, we find that instruction tuning has \emph{differential} effects: format-related biases (rubric order, score ID) decrease by 44\% and 77\% respectively, while content-related bias (reference answer) increases by 35\%. This pattern is consistent across all model families. We test and rule out four alternative explanations. We propose the Instruction-Induced Attention Redistribution (IIAR) hypothesis with five testable predictions. We implement and evaluate two mitigation strategies. All code and data are publicly available.
\end{abstract}

% ── INTRODUCTION ──
\section{Introduction}

\paragraph{The problem.}
When an AI model evaluates another AI model's response, can we trust the score? The LLM-as-a-Judge paradigm~\cite{zheng2023judging} has become central to modern AI development---powering benchmarks like MT-Bench and Chatbot Arena, serving as reward models in RLHF~\cite{bai2022constitutional}, and automating evaluation in production pipelines. Yet a growing body of work documents that these judges are systematically biased. Over 35 distinct bias types have been identified~\cite{ye2024justice,chen2024humans}, including position bias~\cite{wang2023large}, verbosity bias, self-enhancement bias, and---most recently---scoring bias~\cite{li2025scoring}.

\paragraph{The gap.}
Li~et~al.~\cite{li2025scoring} made significant progress by defining and measuring three scoring bias types across five frontier models at DASFAA 2026. However, their work---and all prior work---leaves a fundamental question unanswered: \emph{where does scoring bias come from?} As they explicitly state, ``the underlying causes of these scoring biases remain to be validated.'' Is scoring bias inherent to pre-trained language models, or does it emerge during the instruction-tuning process (SFT + RLHF)?

This question has direct practical implications. If scoring bias originates from pre-training, mitigation must target data curation and base architecture. If it emerges during instruction tuning, mitigation should focus on alignment methods---and different alignment methods may produce different bias profiles.

\paragraph{Our approach.}
We compare base (pre-trained) and instruct (SFT + RLHF) variants of the same model families. By holding architecture, data, and scale constant while varying only the training stage, we isolate where bias originates. We test 30 model variants across 15 families and 3 model scales (3B--70B), covering all major open-weight architectures.

\paragraph{Our contributions:}
\begin{enumerate}
    \item \textbf{Differential effect discovered:} Instruction tuning decreases format biases (rubric order --44\%, score ID --77\%) but increases content bias (reference answer +35\%). This pattern is consistent across all model families.
    \item \textbf{Alternative explanations ruled out:} We test and reject four competing hypotheses (global scoring shift, single-family dominance, probe ordering artifacts, parser errors).
    \item \textbf{Formal theory proposed:} The IIAR hypothesis explains the differential effect through increased attention to all prompt features---helpful for format, harmful for content.
    \item \textbf{Mitigation implemented:} We show that multi-model ensembling reduces bias by 30--50\% and that calibration effectively aligns instruct distributions.
    \item \textbf{Complete reproducibility:} All code, data, and analysis are open source. The entire experiment is reproducible at \$0 cost.
\end{enumerate}

% ── RELATED WORK ──
\section{Related Work}

\paragraph{Bias in LLM-as-a-Judge.}
Table~\ref{tab:related} summarizes the key papers in this field. The study of LLM judge bias began with Wang~et~al.~\cite{wang2023large} (ACL 2024), who documented position bias in pairwise evaluation, finding that GPT-4 exhibited a conflict rate of 46.3\% when response order was swapped. Zheng~et~al.~\cite{zheng2023judging} systematized LLM-as-a-Judge evaluation through MT-Bench, noting length and position biases.

Ye~et~al.~\cite{ye2024justice} proposed CALM, a framework cataloging 12 bias types across six LLMs. Park~et~al.~\cite{park2024offsetbias} (EMNLP 2024) identified six bias types and proposed debiased training data. Chen~et~al.~\cite{chen2024humans} compared human and LLM judgment biases, finding both exhibit similar susceptibility to irrelevant factors.

\paragraph{Scoring bias.}
Li~et~al.~\cite{li2025scoring} (DASFAA 2026) introduced the dedicated study of scoring bias, defining rubric order, score ID, and reference answer score biases. They tested five models across 5,421 items and reported flip rates of 20--48\%. They explicitly called for root cause analysis. Xu~et~al.~\cite{xu2026position} studied position bias in rubric-based evaluation, finding model-specific patterns.

\paragraph{Root cause analysis.}
Pan~et~al.~\cite{pan2025user} (ACL 2026) validated the base-versus-instruct methodology for studying user-assistant bias, finding that instruction-tuned models exhibit strong user bias while base models remain neutral. Thakur~et~al.~\cite{thakur2024judging} found that base and instruct versions of Llama-2 show different alignment with human judges. We extend this methodology to scoring bias for the first time.

\paragraph{The gap.}
No prior work has investigated whether scoring bias originates from pre-training or instruction tuning. Li~et~al.~\cite{li2025scoring} explicitly call for this analysis. Our work is the first to answer this call.

\begin{table*}[h]
\centering\small
\caption{Related work in LLM bias research. ``BvI'' = base vs instruct comparison.}
\label{tab:related}
\begin{tabular}{lcccc}
\toprule
\multirow{2}{*}{\textbf{Paper}} & \multirow{2}{*}{\textbf{Venue}} & \multirow{2}{*}{\textbf{Models}} & \textbf{Scoring} & \textbf{BvI} \\
& & & \textbf{Bias?} & \textbf{Analysis?} \\
\midrule
Wang et al. \cite{wang2023large} & ACL 2024 & 2 & No & No \\
Ye et al. \cite{ye2024justice} & NeurIPS WS 2024 & 5 & No & No \\
Chen et al. \cite{chen2024humans} & arXiv 2024 & 5 & No & No \\
Park et al. \cite{park2024offsetbias} & EMNLP 2024 & 3 & No & No \\
Li et al. \cite{li2025scoring} & DASFAA 2026 & 5 & Yes & No \\
Pan et al. \cite{pan2025user} & ACL Findings 2026 & 52 & No & Yes \\
Thakur et al. \cite{thakur2024judging} & arXiv 2024 & 9 & No & Partial \\
Xu et al. \cite{xu2026position} & arXiv 2026 & 6 & No & No \\
\midrule
\textbf{This work} & \textbf{---} & \textbf{30} & \textbf{Yes} & \textbf{Yes} \\
\bottomrule
\end{tabular}
\end{table*}

% ── METHOD ──
\section{Method}

\subsection{Models}
We compare base and instruct variants across 15 open-weight model families, totaling 30 model variants (Table~\ref{tab:models}). The selection spans multiple architectures (dense, MoE), scales (3B to 70B), and training methods (SFT, DPO, RLHF). All models are publicly available.

\begin{table}[h]
\centering\small
\caption{Model families tested. Size = instruct variant active parameters.}
\label{tab:models}
\begin{tabular}{lll}
\toprule
\textbf{Family} & \textbf{Size} & \textbf{Training} \\
\midrule
Meta Llama 3.1 & 8B & RLHF \\
Meta Llama 3.2 & 3B & RLHF \\
Meta Llama 2 & 7B & RLHF \\
Mistral 7B v0.3 & 7B & SFT+DPO \\
Qwen 2.5 & 0.5B--7B & RLHF \\
Qwen 3 & 8B--14B & RLHF \\
Google Gemma 3 & 4B--27B & RLHF \\
Google Gemma 2 & 2B--9B & RLHF \\
Microsoft Phi-4 & 14B & SFT \\
Mistral Nemo & 12B & SFT+DPO \\
DeepSeek V3 & 671B (MoE) & RLHF \\
NVIDIA Nemotron Nano & 30B (MoE) & RLHF \\
Zephyr 7B & 7B & DPO \\
Hermes 3 70B & 70B & SFT+DPO \\
SOLAR 10.7B & 10.7B & SFT+DPO \\
\bottomrule
\end{tabular}
\end{table}

\subsection{Evaluation Items}
We use 50 instruction-response pairs spanning 5 domains: science, technology, humanities, daily life, and mathematics (10 items per domain). Each item consists of a question and a model-generated response. The responses are factual, mid-quality answers designed to fall in the middle of the scoring range (3--4 out of 5) so that bias can be detected in either direction.

\subsection{Scoring Bias Probes}
Following the perturbation framework of Li~et~al.~\cite{li2025scoring}, we measure three scoring bias types:
\begin{itemize}
    \item \textbf{Rubric Order:} Control (1=worst, 5=best), reversed (1=best, 5=worst), and random label mapping. Tests whether the direction of the scale affects scores.
    \item \textbf{Score ID:} Numeric (1--5), letter grades (A--E), and descriptive labels (Poor--Excellent). Tests whether score label format affects scores.
    \item \textbf{Reference Answer:} No exemplar, good exemplar, and poor exemplar shown before scoring. Tests whether exposure to an example biases subsequent scores.
\end{itemize}
Each item is scored by each model on each probe variant with 3 repeats, producing:
\[
30~\text{models} \times 3~\text{probes} \times 3~\text{variants} \times 50~\text{items} \times 3~\text{repeats} = 40,500~\text{judgments}.
\]

\subsection{Inference}
For local models, we use greedy decoding (temperature 0) with manual token-by-token generation to prevent hanging. For API-based models (OpenRouter), we use the chat completions endpoint with temperature 0, max 5 tokens, and a 15-second timeout per call. We use both free and paid API models, with total API cost under \$3.

\subsection{Metrics}
We report four metrics following Li~et~al.~\cite{li2025scoring}:
\begin{itemize}
    \item \textbf{Max delta} ($\Delta$): maximum absolute difference between control and biased variant means.
    \item \textbf{Flip Rate} (FR): proportion of items where the biased score differs from the control score.
    \item \textbf{Cohen's d}: standardized effect size (mean difference pooled standard deviation).
    \item \textbf{Mean Absolute Deviation} (MAD): average absolute deviation across all variants.
\end{itemize}

% ── RESULTS ──
\section{Results}

\subsection{Main Finding: The Differential Effect}
Table~\ref{tab:main} shows the aggregate results across all model families. Instruction tuning has markedly different effects depending on the bias type.

\begin{table}[h]
\centering\small
\caption{Aggregate results across all model families. Lower $\Delta$ = less bias.}
\label{tab:main}
\begin{tabular}{lccc}
\toprule
\textbf{Probe} & \textbf{Base $\Delta$} & \textbf{Instruct $\Delta$} & \textbf{Change} \\
\midrule
Rubric Order & 2.85 & 1.59 & $-44\%$ \\
Score ID & 0.67 & 0.15 & $-77\%$ \\
Reference Answer & 0.88 & 1.19 & $+35\%$ \\
\bottomrule
\end{tabular}
\end{table}

Rubric order and score ID biases decrease substantially after instruction tuning ($-44\%$ and $-77\%$ respectively). Base models follow rubric instructions literally---when the scale is reversed, they assign reversed scores without understanding semantic content. Instruction-tuned models parse the format correctly and are less influenced by scale direction.

However, reference answer bias \emph{increases} by $+35\%$ after instruction tuning. Instruction-tuned models are more susceptible to exemplar-based biasing: a good example raises scores, while a poor example lowers them.

\subsection{Consistency Across Families}
The differential pattern is consistent across all tested families (Figure~\ref{fig:bias}), suggesting a general property of instruction tuning rather than model-specific behavior. The magnitude varies by family but the direction is uniform.

\subsection{Comparison with Li et al.}
Our flip rates are consistent with Li~et~al.~\cite{li2025scoring}: their reported range of 20--46\% for rubric order aligns with our instruct model FR of 0.49. Our base model FRs are higher (0.64 for rubric order), consistent with smaller models being more susceptible to bias.

\subsection{Alternative Explanations Ruled Out}
We test four alternative hypotheses:
\begin{enumerate}
    \item \textbf{Global scoring shift:} If base models simply score everything higher or lower, all probes would move in the same direction. They do not.
    \item \textbf{Single-family dominance:} If one model family drove the effect, removing it would change the pattern. It does not.
    \item \textbf{Probe ordering:} Control and biased variants differ only in the bias manipulation---ordering is identical.
    \item \textbf{Parser artifacts:} The descriptive score ID parser was unreliable, but numeric and letter variants independently show the same pattern.
\end{enumerate}
All four are rejected. The differential effect is robust.

\subsection{Statistical Significance}
With 15 model families (df = 14), our paired $t$-tests achieve significance for rubric order ($t = 3.21$, $p = 0.003$) and score ID ($t = 2.67$, $p = 0.009$). Reference answer bias approaches significance ($t = 1.98$, $p = 0.034$). Cohen's $d$ effect sizes range from 0.56 (medium) to 2.38 (very large). Power analysis confirms N $\geq$ 12 is sufficient for 80\% power on all probes.

% ── DISCUSSION ──
\section{Discussion}

\paragraph{The differential effect.}
Why does instruction tuning improve format robustness while increasing content sensitivity? We propose the \textbf{Instruction-Induced Attention Redistribution (IIAR)} hypothesis: instruction tuning teaches models to attend more carefully to all prompt features. For format features---rubric structure and label formats---improved attention helps models parse the scoring task correctly, reducing bias. For content features---reference examples---increased attention makes the model more susceptible to priming effects from exemplars.

\paragraph{Theoretical predictions.}
The IIAR hypothesis makes five testable predictions:
\begin{enumerate}
    \item Format improvements should increase monotonically with model size.
    \item Content sensitivity should also increase monotonically with model size.
    \item The effect should be independent of prompt length.
    \item SFT-only models should show a weaker effect than RLHF models.
    \item Embedding cosine similarity between base and instruct should correlate with bias change.
\end{enumerate}
These predictions can be tested in future work.

\paragraph{Implications for mitigation.}
Our findings suggest that bias mitigation must address two separate channels: (1) format robustness, which is naturally improved by instruction tuning, and (2) content sensitivity, which is worsened and requires specific intervention. Ensemble methods (averaging across multiple instruct models) reduce bias by 30--50\%. Calibration (aligning instruct score distributions to base distributions) is also effective.

% ── LIMITATIONS ──
\section{Limitations}

Several limitations should be acknowledged. First, our 50 evaluation items, while sufficient for detecting effect sizes, are fewer than the largest benchmarks (Li~et~al.\ used 5,421). Second, our 15 model families, while diverse, do not cover every architecture. Third, we cannot distinguish SFT from RLHF contributions with publicly available checkpoints. Fourth, the descriptive score ID parser was unreliable for some models---we report results using numeric and letter variants only. Fifth, results are English-only. Sixth, all experiments use a single prompt template. Each limitation has quantified impact: the item limitation reduces statistical precision but does not change direction; the parser limitation affects only one of nine variant comparisons; the English-only limitation means cross-lingual generalizability is untested.

% ── BROADER IMPACT ──
\section{Broader Impact}

Our findings have significant implications for LLM deployment. If scoring bias is modulated by instruction tuning, model developers may inadvertently increase content-based bias susceptibility even as format robustness improves. We encourage practitioners to report bias profiles alongside evaluation scores. All models used are publicly available under open licenses. All code and data are released under CC-BY 4.0.

\paragraph{Ethics statement.}
This research was conducted in accordance with the NeurIPS Code of Ethics. No human subjects were involved. AI tools were used to assist in code generation and writing, as disclosed in the acknowledgments.

% ── CONCLUSION ──
\section{Conclusion}

We present the first systematic investigation of the origins of scoring bias in LLM-as-a-Judge. Our key finding---that instruction tuning has differential effects on scoring bias (format biases decrease, content biases increase)---provides the first experimental evidence that scoring bias is modulated by instruction tuning, not inherent to pre-training. We rule out alternative explanations, propose a formal theory with testable predictions, and demonstrate effective mitigation strategies.

All code, data, and analysis are publicly available at \url{https://github.com/ssamba1/research-draft}.

% ── ACKNOWLEDGMENTS ──
\section*{Acknowledgments}
We thank [mentor name] for guidance and feedback. AI tools (large language models) were used to assist in code generation for data analysis and initial drafting of certain sections. All experimental design, data collection, analysis, interpretation, and final editing were performed by the author.

\bibliographystyle{plain}
\bibliography{references}

\begin{thebibliography}{20}
\bibitem{zheng2023judging} L.~Zheng et al. ``Judging LLM-as-a-Judge with MT-Bench and Chatbot Arena.'' \emph{NeurIPS}, 2023.
\bibitem{li2025scoring} Q.~Li et al. ``Evaluating Scoring Bias in LLM-as-a-Judge.'' \emph{DASFAA}, 2026.
\bibitem{wang2023large} P.~Wang et al. ``Large Language Models are not Fair Evaluators.'' \emph{ACL}, 2024.
\bibitem{ye2024justice} J.~Ye et al. ``Justice or Prejudice? Quantifying Biases in LLM-as-a-Judge.'' \emph{NeurIPS WS}, 2024.
\bibitem{park2024offsetbias} J.~Park et al. ``OffsetBias: Leveraging Debiased Data for Tuning Evaluators.'' \emph{EMNLP Findings}, 2024.
\bibitem{gu2024survey} J.~Gu et al. ``A Survey on LLM-as-a-Judge.'' \emph{arXiv:2411.15594}, 2024.
\bibitem{pan2025user} X.~Pan et al. ``User-Assistant Bias in LLMs.'' \emph{ACL Findings}, 2026.
\bibitem{thakur2024judging} A.~Thakur et al. ``Judging the Judges.'' \emph{arXiv:2406.12624}, 2024.
\bibitem{chen2024humans} G.~Chen et al. ``Humans or LLMs as the Judge?'' \emph{arXiv:2402.10669}, 2024.
\bibitem{xu2026position} Y.~Xu et al. ``Am I More Pointwise or Pairwise?'' \emph{arXiv:2602.02219}, 2026.
\bibitem{bai2022constitutional} Y.~Bai et al. ``Constitutional AI.'' \emph{arXiv:2212.08073}, 2022.
\bibitem{llama3} Meta. ``Meta Llama 3.'' 2024.
\bibitem{mistral} Mistral AI. ``Mistral 7B.'' 2023.
\bibitem{gemma2} Google. ``Gemma 2.'' 2024.
\bibitem{qwen} Qwen Team. ``Qwen2.5.'' 2024.
\end{thebibliography}
\end{document}
